% SOURCE_AUTHORITY: sga5_fr_workpass.tex, lines 14682--14703 (LF-normalized unit).
% SOURCE_SHA256: 791F4EFFC5E02832D5D77ED03518C8156D6F07E4C8238B03545DB93D883FBB28
\subsubsection*{n\textsuperscript{o} 2. Efeito sobre a cohomologia}

\begin{statement}{Proposição 2}
Seja $X$ um pré-esquema de característica $p$.
\begin{enumerate}[label=\alph*)]
\item Para todo feixe de conjuntos $F$ sobre $X_{\et}$, o endomorfismo de $H^0(X,F)$ definido por $(\fr_X,\Fr^*_{F/X})$, ou seja, o composto
\[
H^0(X,F)\longrightarrow H^0(X,\fr_X^*(F))\longrightarrow H^0(X,F),
\]
onde a primeira flecha é a aplicação canônica (caráter contravariante de $H^0(X,F)$ em relação a $X$) e a segunda, $H^0(X,\Fr^*_{F/X})$, é a identidade de $H^0(X,F)$.
\item Para todo feixe de grupos $F$ sobre $X_{\et}$, o endomorfismo de $H^1(X,F)$ definido por $(\fr_X,\Fr^*_{F/X})$, ou seja,
\[
H^1(X,F)\longrightarrow H^1(X,\fr_X^*F)\longrightarrow H^1(X,F),
\]
É a identidade.
\item Para todo complexo $F$ de $\Lambda$-módulos limitado à esquerda sobre $X_{\et}$, o endomorfismo de $\RGamma_X(F)$ definido por $(\fr_X,\Fr^*_{F/X})$, ou seja, o composto
\[
\RGamma_X(F)\longrightarrow \RGamma_X(\fr_X^*F)\longrightarrow \RGamma_X(F),
\]
É a identidade.
\end{enumerate}
\end{statement}
